% Lösungen Einheit 3 von 3 – Wurzelgesetze und rationale Exponenten (zusammenbau v0.1)
\einheitenkopf{Einheit 3 von 3 · Wurzelgesetze und rationale Exponenten}

\erg{17}{a) Produkt oder Quotient: trennbar \quad b) $\sqrt{576} = 24$ und $3 \cdot 8 = 24$ – gleich \quad c) $\frac{8}{3}$ \quad d) $\sqrt{4 \cdot 7} = 2\sqrt{7} \approx 5{,}29$ \quad e) $6\sqrt{3} \approx 10{,}39$ \quad f) $\sqrt{225} = 15$ \quad g) $\sqrt{81} = 9$ \quad h) $4$ \quad i) $5^{\frac{1}{2}+\frac{1}{2}} = 5^1 = 5$ \quad j) $\frac{\sqrt{3}}{3} \approx 0{,}58$ \quad k) $a^{\frac{1}{2}} \cdot b^{\frac{1}{2}} = (a \cdot b)^{\frac{1}{2}}$ – gleicher Exponent, die Basen werden malgenommen. \quad l) $x = \sqrt[3]{512} = 8$ \quad m) $3\sqrt{2} + 2\sqrt{2} = 5\sqrt{2} = 5 \cdot 2^{\frac{1}{2}}$}
\erg{18}{a) Wurzeln aus Summanden darf man nicht zu einer Wurzel zusammenziehen: $\sqrt{81} + \sqrt{144} = 9 + 12 = 21$. \quad b) Etwa $a = 81$, $b = 144$: $\sqrt{225} = 15$, aber $9 + 12 = 21$. \quad c) $x^{\frac{3}{5}}$ \quad d) $a = \sqrt{400 \cdot 2}\,\text{m} = 20\sqrt{2}\,\text{m} \approx 28{,}3\,\text{m}$}
